% DJL: commented out the following... we don't want cite replaced with citep
% default: cite with (Name, year)
%\renewcommand{\cite}{\citep}

\usepackage{longtable}
%%
%% some convenient packages
%%
\usepackage{physics}  %% (p)dv for (partial) derivatives
\usepackage{booktabs}
\usepackage[section]{placeins}
%\usepackage{morefloats} % SAR (2020) commented out
\usepackage{float} % SAR (2020) added
\usepackage{subfig} % SAR (2020) added
\usepackage{amssymb}
\usepackage{setspace}
\usepackage{rotating}
\usepackage{graphicx}
\usepackage{booktabs} 
\usepackage{stmaryrd}
\usepackage{times}
\usepackage{bm}
\usepackage{indentfirst}
\usepackage{natbib}
\usepackage{amsmath}
\usepackage{enumerate}
\usepackage{enumitem}
\usepackage{array}
\usepackage{multicol}

\allowdisplaybreaks   %% allows align environment to break page

% common abbreviations
\newcommand{\eg}{{\it e.g.}\xspace}
\newcommand{\ie}{{\it i.e.}\xspace}
\newcommand{\etc}{{\it etc.}\xspace}
\newcommand{\etal}{\emph{et~al}\mbox{.}\xspace}
\newcommand{\vs}{{vs\mbox{.}}\xspace}

% common Math notation
\newcommand{\NAT}[0]{\mathbb{N}\xspace}
\newcommand{\fun}[1]{\mathit{#1}} % typeset as function name
\newcommand{\setsize}[1]{\left| #1 \right|}
\newcommand{\setdef}[2]{\left\{ #1 \ \left|\  #2\right.\right\}}
\newcommand{\dispsum}[0]{\displaystyle\sum}

\newcommand{\defeq}[0]{\triangleq}
\renewcommand{\mod}{\operatorname{mod}}

%% Macros for literature review
\newcommand*{\Gam}{\mathop{\mathrm{Gamma}}\nolimits}
\newcommand*{\Exp}{\mathop{\mathrm{Exponential}}\nolimits}
\newcommand*{\Ber}{\mathop{\mathrm{Bernoulli}}\nolimits}

% cross referencing equations 
\usepackage{cleveref} % For intelligent cross-referencing
% Define the format for equation references
\crefformat{equation}{(#2#1#3)}

%% Macros for paper 1
\newcommand*{\diag}{\mathop{\mathrm{diag}}\nolimits}
\newcommand*{\bernoulli}{\mathop{\mathrm{Bernoulli}}\nolimits}
\newcommand*{\N}{\mathop{\mathrm{N}}\nolimits}
\newcommand*{\G}{\mathop{\mathrm{Gamma}}\nolimits}
\newcommand*{\W}{\mathop{\mathrm{W}}\nolimits}
\newcommand*{\IW}{\mathop{\mathrm{IW}}\nolimits}
\newcommand*{\TN}{\mathop{\mathrm{TN}}\nolimits}
\newcommand*{\E}{\mathop{\mathrm{E}}\nolimits}
\newcommand*{\Var}{\mathop{\mathrm{Var}}\nolimits}
\newcommand*{\Cov}{\mathop{\mathrm{Cov}}\nolimits}
\newcommand*{\Cor}{\mathop{\mathrm{Cor}}\nolimits}
\newcommand*{\Corr}{\mathop{\mathrm{Corr}}\nolimits}
\newcommand*{\blkdiag}{\mathop{\mathrm{blkdiag}}\nolimits}
\newcommand*{\vect}{\mathop{\mathrm{vec}}\nolimits}
\newcommand*{\POS}{\mathop{\mathrm{POS}}\nolimits}

%% Macros for paper 2
\newcommand*{\CI}{\pi_{\text{CI}}}
\newcommand*{\APP}{\pi_{\text{APP}}}
\newcommand*{\LKJ}{\pi_{\text{LKJ}}}
\newcommand*{\HC}{\pi_{\text{HC}}}
\newcommand*{\CIHC}{\pi_{\text{CIHC}}}
\newcommand*{\pos}{\text{POS}}

%% Macros for paper 3
\def\R{{\mathbb{R}}}
\DeclareMathOperator{\logit}{logit}
\DeclareMathOperator{\expit}{expit}
\DeclareMathOperator{\DY}{\pi_{\text{DY}}}
\DeclareMathOperator{\HPP}{\pi_{\text{HPP}}}
\newcommand{\appropto}{\mathrel{\vcenter{
  \offinterlineskip\halign{\hfil$##$\cr
    \propto\cr\noalign{\kern2pt}\sim\cr\noalign{\kern-2pt}}}}}






% references
\newcommand{\chref}[1]{Chapter~\ref{ch:#1}\xspace}
\newcommand{\chrefs}[2]{Chapters~\ref{ch:#1} and~\ref{ch:#2}\xspace}
\newcommand{\secref}[1]{Section~\ref{sec:#1}\xspace}
\newcommand{\figref}[1]{Figure~\ref{fig:#1}\xspace}
\newcommand{\figrefi}[2]{Figure~\ref{fig:#1}(#2)\xspace}
\newcommand{\tabref}[1]{Table~\ref{tab:#1}\xspace}
\newcommand{\lemref}[1]{Lemma~\ref{lem:#1}\xspace}
\newcommand{\thmref}[1]{Theorem~\ref{thm:#1}\xspace}
\newcommand{\defref}[1]{Definition~\ref{def:#1}\xspace}
\newcommand{\exref}[1]{Example~\ref{ex:#1}\xspace}
\newcommand{\equref}[1]{Equation~(\ref{eq:#1})\xspace}
\newcommand{\inequref}[1]{Inequality~(\ref{eq:#1})\xspace}
\newcommand{\lstref}[1]{Listing~\ref{lst:#1}\xspace}
\newcommand{\pref}[1]{page~\pageref{p:#1}\xspace}
% citations


% special footnotes

% from http://help-csli.stanford.edu/tex/latex-footnotes.shtml
\long\def\symbolfootnote[#1]#2{\begingroup%
\def\thefootnote{\fnsymbol{footnote}}\footnote[#1]{#2}\endgroup}

% Theorems, etc.

\newtheoremstyle{mylemthm}% hnamei 
        {6pt}% hSpace abovei 
        {3pt}% hSpace belowi 
        {\slshape}% hBody fonti 
        {}% hIndent amounti1
        {\bfseries}% hTheorem head fonti 
        {.}% hPunctuation after theorem headi 
        {.5em}% hSpace after theorem headi2
        {}% hTheorem head spec (can be left empty, meaning `normal')i

\theoremstyle{mylemthm}
\newtheorem{assumption}{Assumption}[chapter]
\newtheorem{condition}{Condition}[chapter]

%\theoremstyle{definition}

%% "inherent" from xxexample, but place box at the end of example.
\newenvironment{example}{
\begin{xxexample}
}{
\rightparend{$\Diamond$}
\end{xxexample}
}
% \qed   \sqbullet \blackdiamond \vartriangleleft


% DJL: I added more macros below 

% I think just these ones should suffice
\newcommand{\bs}[1]{\boldsymbol{#1}}
\newcommand{\mc}[1]{\mathcal{#1}}
\newcommand{\overbar}[1]{\mkern 1.5mu\overline{\mkern-1.5mu#1\mkern-1.5mu}\mkern 1.5mu}

% you'll have to check that everything is numbered correctly
\newtheorem{theorem}{Theorem}[chapter]
\newtheorem{lemma}{Lemma}[chapter]
\newtheorem{prop}{Proposition}[chapter]
\newtheorem{corollary}{Corollary}[chapter]
\newtheorem{remark}{Remark}[chapter]
\newtheorem{definition}{Definition}[chapter]

% here's one more
\newcommand{\T}{\intercal}

% and one more
\newcommand\independent{\protect\mathpalette{\protect\independenT}{\perp}}
\def\independenT#1#2{\mathrel{\rlap{$#1#2$}\mkern2mu{#1#2}}}

% last one for real
\newcommand{\pn}{\mathbb{E}_n}

% and one more
\newcommand{\rp}{\mathbb{R}^p}

% added by MTL (2019)
\DeclareMathOperator*{\argmin}{arg\,min}
\DeclareMathOperator*{\argmax}{arg\,max}
\DeclareMathOperator{\imp}{imp}
\DeclareMathOperator{\opt}{opt}
\DeclareMathOperator{\sign}{sign}
\newcommand{\bignorm}[1]{\left\lVert#1\right\rVert}
\newcommand{\bbR}{\mathbb{R}}
\newcommand{\bbE}{\mathbb{E}}
\newcommand{\calX}{\mathcal{X}}
\newcommand{\calY}{\mathcal{Y}}
\newcommand{\calP}{\mathcal{P}}
\newcommand{\calN}{\mathcal{N}}
\newcommand{\meani}{\displaystyle{\frac{1}{n}\sum_{i=1}^{n}}}
\newcommand{\sumi}{\displaystyle{\sum_{i=1}^{n}}}
\newcommand{\bhat}{\hat{\beta}}
\newcommand{\btil}{\tilde{\beta}}
\newcommand{\Dhat}{\hat{\Delta}}
\newcommand{\Szsig}{S^{0}_{\sigma}}

%% Added by SAR (2020):

% for wide hat
\usepackage{scalerel,stackengine}
\stackMath
\newcommand\reallywidehat[1]{%
\savestack{\tmpbox}{\stretchto{%
  \scaleto{%
    \scalerel*[\widthof{\ensuremath{#1}}]{\kern.1pt\mathchar"0362\kern.1pt}%
    {\rule{0ex}{\textheight}}%WIDTH-LIMITED CIRCUMFLEX
  }{\textheight}% 
}{2.4ex}}%
\stackon[-6.9pt]{#1}{\tmpbox}%
}

% for commenting
\newcommand{\sarah}[1]{\textcolor{green}{(Sarah: #1)}}

% to name rows of subfigs
% https://tex.stackexchange.com/questions/239715/add-titles-for-rows-and-columns-in-a-subfloat/239809
\usepackage{mwe}
\newlength{\tempheight}
\newlength{\tempwidth}
\newcommand{\rowname}[1]% #1 = text
{\rotatebox{90}{\makebox[\tempheight][c]{\textbf{#1}}}}
\newcommand{\columnname}[1]% #1 = text
{\makebox[\tempwidth][c]{\textbf{#1}}}

% Begin code from  PDF -> Latex (Rmd generated tex file)
\usepackage{color}
\usepackage{fancyvrb}
\newcommand{\VerbBar}{|}
\newcommand{\VERB}{\Verb[commandchars=\\\{\}]}
\DefineVerbatimEnvironment{Highlighting}{Verbatim}{commandchars=\\\{\}}
% Add ',fontsize=\small' for more characters per line
\usepackage{framed}
\definecolor{shadecolor}{RGB}{248,248,248}
\newenvironment{Shaded}{\begin{snugshade}}{\end{snugshade}}
\newcommand{\KeywordTok}[1]{\textcolor[rgb]{0.13,0.29,0.53}{\textbf{#1}}}
\newcommand{\DataTypeTok}[1]{\textcolor[rgb]{0.13,0.29,0.53}{#1}}
\newcommand{\DecValTok}[1]{\textcolor[rgb]{0.00,0.00,0.81}{#1}}
\newcommand{\BaseNTok}[1]{\textcolor[rgb]{0.00,0.00,0.81}{#1}}
\newcommand{\FloatTok}[1]{\textcolor[rgb]{0.00,0.00,0.81}{#1}}
\newcommand{\ConstantTok}[1]{\textcolor[rgb]{0.00,0.00,0.00}{#1}}
\newcommand{\CharTok}[1]{\textcolor[rgb]{0.31,0.60,0.02}{#1}}
\newcommand{\SpecialCharTok}[1]{\textcolor[rgb]{0.00,0.00,0.00}{#1}}
\newcommand{\StringTok}[1]{\textcolor[rgb]{0.31,0.60,0.02}{#1}}
\newcommand{\VerbatimStringTok}[1]{\textcolor[rgb]{0.31,0.60,0.02}{#1}}
\newcommand{\SpecialStringTok}[1]{\textcolor[rgb]{0.31,0.60,0.02}{#1}}
\newcommand{\ImportTok}[1]{#1}
\newcommand{\CommentTok}[1]{\textcolor[rgb]{0.56,0.35,0.01}{\textit{#1}}}
\newcommand{\DocumentationTok}[1]{\textcolor[rgb]{0.56,0.35,0.01}{\textbf{\textit{#1}}}}
\newcommand{\AnnotationTok}[1]{\textcolor[rgb]{0.56,0.35,0.01}{\textbf{\textit{#1}}}}
\newcommand{\CommentVarTok}[1]{\textcolor[rgb]{0.56,0.35,0.01}{\textbf{\textit{#1}}}}
\newcommand{\OtherTok}[1]{\textcolor[rgb]{0.56,0.35,0.01}{#1}}
\newcommand{\FunctionTok}[1]{\textcolor[rgb]{0.00,0.00,0.00}{#1}}
\newcommand{\VariableTok}[1]{\textcolor[rgb]{0.00,0.00,0.00}{#1}}
\newcommand{\ControlFlowTok}[1]{\textcolor[rgb]{0.13,0.29,0.53}{\textbf{#1}}}
\newcommand{\OperatorTok}[1]{\textcolor[rgb]{0.81,0.36,0.00}{\textbf{#1}}}
\newcommand{\BuiltInTok}[1]{#1}
\newcommand{\ExtensionTok}[1]{#1}
\newcommand{\PreprocessorTok}[1]{\textcolor[rgb]{0.56,0.35,0.01}{\textit{#1}}}
\newcommand{\AttributeTok}[1]{\textcolor[rgb]{0.77,0.63,0.00}{#1}}
\newcommand{\RegionMarkerTok}[1]{#1}
\newcommand{\InformationTok}[1]{\textcolor[rgb]{0.56,0.35,0.01}{\textbf{\textit{#1}}}}
\newcommand{\WarningTok}[1]{\textcolor[rgb]{0.56,0.35,0.01}{\textbf{\textit{#1}}}}
\newcommand{\AlertTok}[1]{\textcolor[rgb]{0.94,0.16,0.16}{#1}}
\newcommand{\ErrorTok}[1]{\textcolor[rgb]{0.64,0.00,0.00}{\textbf{#1}}}
\newcommand{\NormalTok}[1]{#1}
% End code from  PDF -> Latex (Rmd generated tex file)

% Personal project commands
\input{0.header_thesis}